% ============================================================
% PROMPT 1: In-Group / Out-Group Rhetoric
% ============================================================

\begin{tcolorbox}[
  title=Prompt 1: In-Group / Out-Group Rhetoric Scoring,
  fonttitle=\bfseries,
  colback=gray!5,
  colframe=black!60,
  breakable
]

You are an expert in political rhetoric and social identity theory.
Your task is to analyze a single line of dialogue from the anime ``Attack on Titan''
and score it on three indicators of in-group / out-group rhetoric:
\begin{enumerate}
  \item Boundary Marking
  \item Threat Framing
  \item Solidarity Appeal
\end{enumerate}

For each indicator, apply the step-by-step reasoning process described below,
then assign a score on the following scale:

\medskip
\textbf{Scoring Scale:}
\begin{itemize}
  \item 1 = Absent --- no evidence of this indicator in the text
  \item 2 = Very weak --- barely perceptible, highly implicit or ambiguous
  \item 3 = Moderate --- present but partial; a reasonable reader might debate it
  \item 4 = Strong --- clear and explicit use of this indicator
  \item 5 = Very strong --- central, unmistakable, defining the statement
\end{itemize}

\medskip\hrule\medskip

\textbf{Indicator 1: Boundary Marking}

\textit{Definition:} The explicit or implicit linguistic construction of an us-versus-them
distinction through group labels, pronouns, or categorical separation between
in-group and out-group members.

\textit{Reasoning Process:}\\
Step 1: Identify any group labels or categorical terms (e.g., Eldians, Marleyans,
devils, us, them, our people).\\
Step 2: Determine whether the sentence draws a clear line between ``we/us'' and
``they/them'' as distinct social categories.\\
Step 3: Assess whether the distinction is merely descriptive or identity-constitutive.\\
Step 4: Assign a score 1--5.

\textit{Few-Shot Examples:}

\begin{quote}
\textbf{Example 1}\\
Text: ``We Eldians have been persecuted by Marleyans for generations. That is simply what they do to our kind.''\\
Step 1: Group labels --- ``Eldians'', ``Marleyans'', ``our kind''.\\
Step 2: Explicit we/they distinction.\\
Step 3: Identity-constitutive.\\
\textbf{Score: 5}
\end{quote}

\begin{quote}
\textbf{Example 2}\\
Text: ``Those people beyond the walls have always looked down on us.''\\
Step 1: ``Those people beyond the walls'' vs.\ implicit ``us''.\\
Step 2: Clear we/they distinction.\\
Step 3: Identity-constitutive --- ``always'' frames persistent group-based contempt.\\
\textbf{Score: 4}
\end{quote}

\begin{quote}
\textbf{Example 3}\\
Text: ``People like us don't get to choose how we live.''\\
Step 1: ``People like us'' implies a social category.\\
Step 2: Implicit we/they distinction.\\
Step 3: Ambiguous --- constructs an in-group identity but does not explicitly oppose or name an out-group.\\
\textbf{Score: 3}
\end{quote}

\begin{quote}
\textbf{Example 4}\\
Text: ``The enemy soldiers have breached the wall on the eastern side.''\\
Step 1: ``Enemy soldiers'' present, but no specific group identity label.\\
Step 2: No corresponding in-group identity invoked.\\
Step 3: Descriptive tactical statement, not an identity claim.\\
\textbf{Score: 1}
\end{quote}

\medskip\hrule\medskip

\textbf{Indicator 2: Threat Framing}

\textit{Definition:} The portrayal of an out-group as posing an existential, physical,
moral, or civilizational threat to the in-group's survival, safety, or way of life.

\textit{Reasoning Process:}\\
Step 1: Identify whether an out-group is mentioned or implied.\\
Step 2: Determine whether the out-group is characterized as dangerous or harmful.\\
Step 3: Assess severity --- mundane conflict vs.\ existential threat.\\
Step 4: Assign a score 1--5.

\textit{Few-Shot Examples:}

\begin{quote}
\textbf{Example 1}\\
Text: ``As long as those people exist beyond the walls, none of us will ever be truly safe. They will not stop until we are gone.''\\
Step 1: Out-group implied.\\
Step 2: Persistent, intentional threat framed.\\
Step 3: Existential --- ``until we are gone'' is eliminationist.\\
\textbf{Score: 5}
\end{quote}

\begin{quote}
\textbf{Example 2}\\
Text: ``Marley will come for us eventually. That's just how empires work.''\\
Step 1: Out-group explicitly named.\\
Step 2: Future threat implied.\\
Step 3: Moderately existential.\\
\textbf{Score: 4}
\end{quote}

\begin{quote}
\textbf{Example 3}\\
Text: ``They've been stirring up trouble in the outer districts for years.''\\
Step 1: Out-group implied.\\
Step 2: Persistent but low-level problem.\\
Step 3: Not existential.\\
\textbf{Score: 3}
\end{quote}

\begin{quote}
\textbf{Example 4}\\
Text: ``The scouts have reported movement near the southern border.''\\
Step 1: No out-group explicitly characterized.\\
Step 3: Neutral tactical observation.\\
\textbf{Score: 1}
\end{quote}

\medskip\hrule\medskip

\textbf{Indicator 3: Solidarity Appeal}

\textit{Definition:} An invocation of shared fate, collective sacrifice, common mission,
or mutual obligation among in-group members, used to reinforce group cohesion
and motivate collective action.

\textit{Reasoning Process:}\\
Step 1: Identify whether the sentence addresses a collective in-group.\\
Step 2: Determine whether the sentence invokes shared experience --- suffering together, fighting together, or being bound by the same fate.\\
Step 3: Assess whether the appeal is directed at motivating or emotionally binding the in-group.\\
Step 4: Assign a score 1--5.

\textit{Few-Shot Examples:}

\begin{quote}
\textbf{Example 1}\\
Text: ``We have all lost something. Every one of us carries that weight. But that is exactly why we cannot stop fighting --- for those who are no longer here.''\\
Step 1: Collective in-group --- ``we'', ``every one of us''.\\
Step 2: Shared suffering invoked.\\
Step 3: Explicitly motivational.\\
\textbf{Score: 5}
\end{quote}

\begin{quote}
\textbf{Example 2}\\
Text: ``If we're going to die anyway, let's at least die together as soldiers.''\\
Step 1: Collective in-group --- ``we'', ``soldiers''.\\
Step 2: Shared fate invoked.\\
Step 3: Solidarity appeal even if the emotional register is resigned.\\
\textbf{Score: 4}
\end{quote}

\begin{quote}
\textbf{Example 3}\\
Text: ``We need to look out for each other out there.''\\
Step 1: Collective in-group --- ``we'', ``each other''.\\
Step 2: Mutual obligation invoked.\\
Step 3: Moderately motivational.\\
\textbf{Score: 3}
\end{quote}

\begin{quote}
\textbf{Example 4}\\
Text: ``I joined the Survey Corps three years ago.''\\
Step 1: References in-group membership only biographically.\\
Step 2: No shared fate or collective emotional bond.\\
Step 3: Descriptive, not motivational.\\
\textbf{Score: 1}
\end{quote}

\medskip\hrule\medskip

\textbf{Output Format}

Return output strictly as JSON. Do not include any text outside the JSON block.

\begin{verbatim}
{
  "analysis": [
    {"strategy": "Boundary Marking",  "score": <1|2|3|4|5>},
    {"strategy": "Threat Framing",    "score": <1|2|3|4|5>},
    {"strategy": "Solidarity Appeal", "score": <1|2|3|4|5>}
  ]
}
\end{verbatim}

\end{tcolorbox}


% ============================================================
% PROMPT 2: Moral Disengagement
% ============================================================

\begin{tcolorbox}[
  title=Prompt 2: Moral Disengagement Strategy Scoring,
  fonttitle=\bfseries,
  colback=gray!5,
  colframe=black!60,
  breakable
]

You are a political psychology text-analysis assistant trained to identify mechanisms
of moral disengagement in narrative, political, and conflict-related texts.

Your task is to evaluate whether the text contains any of the following eight moral
disengagement strategies:
\begin{enumerate}
  \item Moral Justification
  \item Euphemistic Labeling
  \item Advantageous Comparison
  \item Displacement of Responsibility
  \item Diffusion of Responsibility
  \item Disregard / Distortion of Consequences
  \item Dehumanization
  \item Attribution of Blame
\end{enumerate}

\medskip\hrule\medskip

\textbf{General Analysis Instructions}

For each strategy:
\begin{itemize}
  \item Carefully reason through the meaning of the text step by step before assigning a score.
  \item Focus on implied moral framing, justification, and portrayal of victims or responsibility.
  \item Consider both explicit and implicit language.
  \item Do not rely solely on keywords.
  \item Use the examples as guidance but not as strict templates.
\end{itemize}

\textit{Scoring Rules:}
\begin{itemize}
  \item 1 = Absent --- no evidence of this strategy
  \item 2 = Very weak --- barely perceptible, highly implicit or ambiguous
  \item 3 = Moderate --- present but secondary; a reasonable reader might debate it
  \item 4 = Strong --- clear and explicit use of this strategy
  \item 5 = Very strong --- central, unmistakable; defining the moral framing
\end{itemize}

\textit{Pre-Scoring Checklist (apply before every evaluation):}
\begin{enumerate}
  \item Identify whether harmful conduct is being discussed.
  \item Determine whether the text attempts to morally justify, normalize, minimize, or excuse such conduct.
  \item Evaluate which specific mechanism is being used.
\end{enumerate}

\medskip\hrule\medskip

\textbf{1. Moral Justification}

\textit{Definition:} Harmful or violent actions are framed as necessary for a higher moral
purpose such as survival, justice, freedom, security, patriotism, revenge, or protection of the group.

\begin{quote}
\textbf{Example 1}\\
Text: ``We must eliminate them to protect the future of our nation.''\\
Reasoning: Frames lethal violence as morally necessary for national survival.\\
\textbf{Score: 5}

\textbf{Example 2}\\
Text: ``War is unfortunate, but sometimes sacrifices are required for peace.''\\
Reasoning: Morally legitimizes harmful actions as necessary sacrifices, but hedges.\\
\textbf{Score: 4}

\textbf{Example 3}\\
Text: ``They attacked us first, so we fought back.''\\
Reasoning: Mild retaliatory framing, not a positive moral mission.\\
\textbf{Score: 2}
\end{quote}

\medskip\hrule\medskip

\textbf{2. Euphemistic Labeling}

\textit{Definition:} Harmful or violent actions are described using sanitized, indirect,
technical, or emotionally softened language.

\begin{quote}
\textbf{Example 1}\\
Text: ``The operation successfully neutralized hostile targets.''\\
Reasoning: ``Neutralized'' replaces ``killed''; ``hostile targets'' replaces ``people.''\\
\textbf{Score: 5}

\textbf{Example 2}\\
Text: ``Collateral damage was unavoidable during the mission.''\\
Reasoning: Established euphemism for civilian death.\\
\textbf{Score: 5}

\textbf{Example 3}\\
Text: ``The army attacked the village.''\\
Reasoning: Direct, unambiguous language. No softening.\\
\textbf{Score: 1}
\end{quote}

\medskip\hrule\medskip

\textbf{3. Advantageous Comparison}

\textit{Definition:} Harmful actions are justified by comparing them to worse actions
committed by others.

\begin{quote}
\textbf{Example 1}\\
Text: ``At least we did not massacre civilians like they did.''\\
Reasoning: Directly minimizes own actions by contrasting with a more extreme atrocity.\\
\textbf{Score: 5}

\textbf{Example 2}\\
Text: ``Our methods may be harsh, but they are far more humane than theirs.''\\
Reasoning: Acknowledges harshness but repositions own conduct as morally superior.\\
\textbf{Score: 4}

\textbf{Example 3}\\
Text: ``They committed terrible crimes.''\\
Reasoning: Describes the other group's actions but involves no comparison with own conduct.\\
\textbf{Score: 1}
\end{quote}

\medskip\hrule\medskip

\textbf{4. Displacement of Responsibility}

\textit{Definition:} Individuals portray harmful actions as the result of orders, authority,
duty, or external pressure rather than personal choice.

\begin{quote}
\textbf{Example 1}\\
Text: ``I was only following orders.''\\
Reasoning: Responsibility unambiguously shifted to the authority issuing orders.\\
\textbf{Score: 5}

\textbf{Example 2}\\
Text: ``The commander required us to carry out the attack.''\\
Reasoning: Frames the harmful act as an external obligation.\\
\textbf{Score: 4}

\textbf{Example 3}\\
Text: ``We decided to attack.''\\
Reasoning: Agency is clearly personal and direct. No displacement.\\
\textbf{Score: 1}
\end{quote}

\medskip\hrule\medskip

\textbf{5. Diffusion of Responsibility}

\textit{Definition:} Responsibility for harmful actions is spread across a group or collective.

\begin{quote}
\textbf{Example 1}\\
Text: ``Everyone supported the decision.''\\
Reasoning: Responsibility distributed across an entire collective.\\
\textbf{Score: 4}

\textbf{Example 2}\\
Text: ``We were all involved in the operation.''\\
Reasoning: Collective involvement dilutes individual accountability.\\
\textbf{Score: 4}

\textbf{Example 3}\\
Text: ``I personally ordered the attack.''\\
Reasoning: Responsibility individualized and directly claimed.\\
\textbf{Score: 1}
\end{quote}

\medskip\hrule\medskip

\textbf{6. Disregard / Distortion of Consequences}

\textit{Definition:} Harmful outcomes are minimized, dismissed, ignored, or portrayed
as insignificant.

\begin{quote}
\textbf{Example 1}\\
Text: ``The civilian damage was minimal.''\\
Reasoning: Directly minimizes harm done to civilians.\\
\textbf{Score: 4}

\textbf{Example 2}\\
Text: ``The victims will recover quickly.''\\
Reasoning: Downplays seriousness by implying swift recovery.\\
\textbf{Score: 3}

\textbf{Example 3}\\
Text: ``Thousands of civilians were killed.''\\
Reasoning: Consequences described directly without minimization.\\
\textbf{Score: 1}
\end{quote}

\medskip\hrule\medskip

\textbf{7. Dehumanization}

\textit{Definition:} Individuals or groups are portrayed as less than human, animal-like,
impure, evil, or fundamentally inferior.

\begin{quote}
\textbf{Example 1}\\
Text: ``They are parasites destroying our society.''\\
Reasoning: ``Parasites'' is a classic dehumanizing metaphor depicting the target as subhuman.\\
\textbf{Score: 5}

\textbf{Example 2}\\
Text: ``Those monsters deserve extinction.''\\
Reasoning: ``Monsters'' denies human qualities; ``deserve extinction'' positions the group outside moral concern.\\
\textbf{Score: 5}

\textbf{Example 3}\\
Text: ``Our enemies are dangerous.''\\
Reasoning: Describes a threat but does not strip targets of human status.\\
\textbf{Score: 1}
\end{quote}

\medskip\hrule\medskip

\textbf{8. Attribution of Blame}

\textit{Definition:} Victims or targeted groups are portrayed as responsible for
provoking or deserving the harm directed against them.

\begin{quote}
\textbf{Example 1}\\
Text: ``They brought this destruction upon themselves.''\\
Reasoning: Victims directly and explicitly blamed for the violence directed at them.\\
\textbf{Score: 5}

\textbf{Example 2}\\
Text: ``If they had obeyed us, none of this would have happened.''\\
Reasoning: Harm presented as a consequence of the victims' own noncompliance.\\
\textbf{Score: 5}

\textbf{Example 3}\\
Text: ``The conflict began after protests erupted.''\\
Reasoning: Describes a chronological sequence without assigning moral fault.\\
\textbf{Score: 1}
\end{quote}

\medskip\hrule\medskip

\textbf{Output Format}

Return output strictly as JSON. Do not include any text outside the JSON block.

\begin{verbatim}
{
  "analysis": [
    {"strategy": "Moral Justification",                   "score": <1|2|3|4|5>},
    {"strategy": "Euphemistic Labeling",                  "score": <1|2|3|4|5>},
    {"strategy": "Advantageous Comparison",               "score": <1|2|3|4|5>},
    {"strategy": "Displacement of Responsibility",        "score": <1|2|3|4|5>},
    {"strategy": "Diffusion of Responsibility",           "score": <1|2|3|4|5>},
    {"strategy": "Disregard / Distortion of Consequences","score": <1|2|3|4|5>},
    {"strategy": "Dehumanization",                        "score": <1|2|3|4|5>},
    {"strategy": "Attribution of Blame",                  "score": <1|2|3|4|5>}
  ]
}
\end{verbatim}

\end{tcolorbox}
